% Generated draft: requires mathematical and source review.
\documentclass[11pt]{article}
\usepackage[margin=25mm]{geometry}
\usepackage{fontspec}
\setmainfont{Latin Modern Roman}
\usepackage{amsmath,amssymb,mathtools,unicode-math}
\setmathfont{Latin Modern Math}
\usepackage{enumitem,needspace,xurl,hyperref}
\setlength{\parindent}{0pt}
\setlength{\parskip}{5pt}
\setlength{\emergencystretch}{3em}
\newcommand{\N}{\mathbb{N}}
\newcommand{\Z}{\mathbb{Z}}
\newcommand{\Q}{\mathbb{Q}}
\newcommand{\R}{\mathbb{R}}
\newcommand{\C}{\mathbb{C}}
\newcommand{\F}{\mathbb{F}}
\newcommand{\A}{\mathbb{A}}
\newcommand{\PP}{\mathbb{P}}
\newcommand{\E}{\mathbb{E}}
\newcommand{\eps}{\varepsilon}
\newcommand{\dd}{\,\mathrm{d}}
\begin{document}
\section*{UnsolvedMath ID: 30001716}
\textbf{Source problem code:} catalogue entry
\subsection{Definitions and mathematical statement}
Let \(\Gamma\) be a quiver and let \(v=(v_i)\) be a dimension vector indexed by its vertices. The supplied formulation describes \(C_{\Gamma,v}\) as a projective curve having one rational component of multiplicity \(v_i\) for each vertex \(i\), with nodal intersections along the edges of \(\Gamma\). It denotes its compactified Jacobian by \(J_v\) and the associated Crawley--Boevey--Shaw character variety by \(\mathcal M_v\). The ground field, the precise curve construction, and the parameters defining these moduli spaces are not supplied; these descriptions therefore do not uniquely specify the objects.

The proposed graphical non-Abelian Hodge correspondence, attributed by the supplied record to Conjecture 5 of the cited report, asks for graded cohomology isomorphisms
\[
\Phi_v:H^*(J_v)\xrightarrow{\sim}H^*(\mathcal M_v)
\]
identifying the indicated filtration pieces:
\[
\Phi_v\bigl(P_kH^*(J_v)\bigr)=W_{2k}H^*(\mathcal M_v).
\]
Here \(W\) denotes the indicated weight filtration on character-variety cohomology, while \(P\) denotes the indicated filtration on compactified-Jacobian cohomology. The supplied record does not specify the cohomology coefficients, the construction and normalization of \(P\), or the range of the filtration index \(k\).

The correspondence is also asserted to yield the polynomial identity
\[
P(H;J_v;t)=P(P;\mathcal M_v;t)
=t^{2d_v}A_\Gamma(v,t^{-2}).
\]
The symbols \(P(H;J_v;t)\), \(P(P;\mathcal M_v;t)\), \(d_v\), and \(A_\Gamma(v,t)\) are retained exactly as supplied: their definitions and grading conventions are absent. In particular, the occurrence of \(P(P;\mathcal M_v;t)\) does not specify how that polynomial relates to the filtration \(P\) previously placed on \(H^*(J_v)\).

The supplied scope is arbitrary quivers and dimension vectors, with no further hypotheses recorded. This is a statement-only transcription of the proposed correspondence; the missing constructions, parameters, and conventions prevent a fully specified mathematical conjecture from being recovered from this record alone.
\subsection{Short English statement}
The supplied conjecture proposes a graphical non-Abelian Hodge correspondence between the cohomology of a compactified Jacobian and a Crawley--Boevey--Shaw character variety. It should match the stated filtrations and give the displayed polynomial identity. Essential geometric choices and polynomial definitions are missing, so the formulation requires human mathematical review. The supplied open status is dated curation, not an independently verified conclusion.
\subsection{Sources}
\begin{itemize}
\item Supplied local UnsolvedMath registry record, permanent ID \(30001716\), title \emph{Graphical Non-Abelian Hodge Correspondence}, source problem code \texttt{catalogue entry}, in the Oberwolfach Reports Open Problems collection. The record gives proposed year 2011 and names no individual proposer.
\item \emph{Representation Theory of Quivers and Finite Dimensional Algebras} (2011), Oberwolfach Reports, pp.\ 523--608, DOI \texttt{10.4171/owr/2011/10}. Supplied source and literature-source link, labeled ``Oberwolfach Report 10/2011'': \url{https://doi.org/10.4171/owr/2011/10}. The supplied assessment identifies the correspondence as Conjecture 5 for arbitrary quivers.
\item Literature assessment date: \texttt{2026-08-21}. The supplied assessment labels the problem open and states that the report and subsequent literature located in that curation do not supply a general proof. This is supplied dated curation, not an independently verified literature assessment; the background expressly requests independent verification of time-sensitive status and cited claims.
\item The registry labels the statement \texttt{corrected\_verified} and reports manual semantic, typographic, and LaTeX review, describing the cleanup as self-contained and meaning-preserving. These are supplied editorial claims; they do not resolve the missing mathematical specifications identified here.
\end{itemize}
\end{document}
